\documentclass[reprint,superscriptaddress,nofootinbib,nobibnotes,amsmath,amssymb,aps,prd,floatfix,longbibliography]{revtex4-2}
\usepackage[T1]{fontenc}
\usepackage[utf8]{inputenc}
\usepackage{amsmath,amssymb,bm,booktabs,graphicx,xcolor}
\usepackage{hyperref}
\usepackage{xurl}
\def\btemubody{body.tex}
\def\btemubib{references}
\graphicspath{{figures/}}
\hypersetup{colorlinks=true,linkcolor=blue,citecolor=blue,urlcolor=blue}
\providecommand{\tightlist}{\setlength{\itemsep}{0pt}\setlength{\parskip}{0pt}}
\setlength{\emergencystretch}{2em}
\renewcommand{\dbltopfraction}{0.95}
\renewcommand{\textfraction}{0.02}
\begin{document}
\title{BTemu: a history emulator for BT-modified reionization calculations}
\author{Haina Huang}
\affiliation{Department of Physics, The Chinese University of Hong Kong, Shatin, Hong Kong, China}
\date{November 6, 2026}
\begin{abstract}
Small-scale matter spectra and galaxy formation both affect reionization, so spectral inference requires marginalization over galaxy properties. BTemu predicts 32-node neutral-fraction histories from ten galaxy, heating and spectrum inputs. We use it to vary six galaxy parameters, a matter-spectrum break scale $k_p$ and slope $m_s$, with the X-ray settings fixed and the luminosity-function (LF) calculation exact. The 100,000-history model reports development trajectory-RMSE median $8.51\times10^{-4}$ and 95th percentile $5.56\times10^{-3}$. In a representative comparison with the power-law (PL) reference, blue-tilted (BT) samples at $k_p=1\,h\,\mathrm{Mpc}^{-1}$ and $m_s=1.5$ select lower star formation per unit halo mass, with a stronger reduction at lower mass. Marginalizing six galaxy coordinates from the last 300 steps of a 600-step joint chain gives 95\% sample intervals $k_p\in[4.420,29.071]\,h\,\mathrm{Mpc}^{-1}$ and $m_s\in[0.584,3.962]$. These preliminary representative-history results show a change at low $k_p$ and little range reduction for $m_s$ relative to the uniform-prior reference.
\end{abstract}
\maketitle
\begin{quote}\small Working draft v0.10. Joint spectrum--galaxy inference.\end{quote}
\input{\btemubody}
\end{document}
